\section{The lowest component and the missing-edge charge}
The component operation used above costs no extra gates in the edge-support class, but does not by itself charge the missing edges. The underlying lower-envelope theorem is classical \citep[Theorem~2.4, p.~879]{JerrumSnir1982}; we spell out its application to the precise bidegree question.

\begin{proposition}\label{prop:componentzero}
Let $G=(X\sqcup Y,E)$ be bipartite, and let $P\in\mathbb N[X,Y]$ have the property that every monomial support contains an edge of $G$. A positive arithmetic circuit for $P$ gives a circuit for its bidegree-$(1,1)$ component with no increase in either additions or multiplications, provided a free zero output is permitted when that component is zero. If the component is nonzero, no constant input need be introduced.
\end{proposition}
\begin{proof}
For a nonzero polynomial $A$, let $\delta(A)$ be its least total degree and $\ell(A)$ its component in that degree. Positivity gives
\[
 \delta(AB)=\delta(A)+\delta(B),\qquad
 \ell(AB)=\ell(A)\ell(B).
\]
For addition, retain only the operand of smaller minimum degree when the two degrees differ; when they agree, add their lower components. Apply these rules in topological order, preserving every shared wire, and delete unused gates. No gate introduces more than one operation of its original type. If $P$ has a quadratic monomial, its minimum degree is two, and every quadratic monomial has one variable from each side. Therefore $\ell(P)=P_{1,1}$. If there is no quadratic monomial, the desired component is zero.
\end{proof}

Write $E'$ for the quadratic support. The extracted circuit computes a positive weighting of $H=(V,E')$, so its lower bound is $|E'|+\tau(H)-1$. To recover the conjectured bound for $G$, a separate argument must charge the deleted gates for both $|E\setminus E'|$ and $\tau(G)-\tau(H)$. A charge involving only additions is already impossible for the Petersen example in \citet[Section~6, example `A repeated-variable tradeoff on the Petersen graph']{OneSidedCompanion}.

Here is the exact computation. Put
\[
 u=x_0+x_6,\quad v=x_5+x_6,\quad w=x_0+x_9,
 \quad t=x_9+x_5,\quad p=x_7+x_3.
\]
The polynomial
\[
 P=x_4(w+x_3)+(x_1u+x_2)(p+x_1)
       +x_8(x_3+v)+(x_7+uvw)t
\]
has precisely the Petersen edges as its minimal monomial supports. Its
component of total degree two is
\[
 Q=x_4(w+x_3)+x_2(p+x_1)+x_8(x_3+v)+x_7t.
\]
This example concerns total degree, since the Petersen graph is not
bipartite. The missing quadratic edges are $01,05,16,69$; the other eleven are
\[
 04,34,49,12,23,27,38,58,68,57,79.
\]
Indeed the terms of $Q$ give those eleven edges. The remaining terms are $x_1u(p+x_1)$ and $uvwt$. Every monomial of the former contains $01$ or $16$, and its $x_1^2u$ terms witness both pairs. For the latter, Boolean absorption gives
\[
 (x_0\vee x_6)(x_5\vee x_6)(x_0\vee x_9)(x_9\vee x_5)
 =(x_0x_5)\vee(x_6x_9).
\]
To check this identity, the first two factors reduce to $x_6\vee x_0x_5$, the last two to $x_9\vee x_0x_5$, and their product is the displayed disjunction. Positivity then shows that every monomial support contains $05$ or $69$, and both occur. Thus the fifteen minimal supports are exactly the asserted graph; its usual outer-five-cycle and inner-star presentation, after relabelling, identifies it as the Petersen graph.


The displayed DAG for $P$ uses thirteen additions and seven products.
Its lower component $Q$ uses ten additions and four products after
unused wires are removed. Four edges have therefore lost quadratic
witnesses while only three additions are deleted. The total loss is six
gates, consistent with a possible joint charge. This example supplies no
counterexample to the total-size conjecture.
